\documentclass[11pt,a4paper]{article}

% ── Packages ──
\usepackage[utf8]{inputenc}
\usepackage[T1]{fontenc}
\usepackage{amsmath,amssymb}
\usepackage{siunitx}
\usepackage{booktabs}
\usepackage{geometry}
\usepackage{hyperref}
\usepackage{float}

\geometry{margin=2.5cm}
\hypersetup{colorlinks=true,linkcolor=blue,citecolor=blue,urlcolor=blue}

\title{MARS Geometries: Propeller vs.\ Rotating Disk}
\author{}
\date{}

\begin{document}
\maketitle
\tableofcontents
\newpage

% ══════════════════════════════════════════════════════════════
\section{Pressure Vessel and Gas Environment}
\label{sec:vessel}
% ══════════════════════════════════════════════════════════════

Both geometries operate inside the same cylindrical pressure vessel.

\paragraph{Vessel.}
The vessel is a vertical cylinder with inner radius
\begin{equation}
  R_\mathrm{cyl} = R + \frac{\Delta D}{2}\,,
\end{equation}
where $R = \SI{1.6}{m}$ is the radius of the rotating element
(propeller or disk) and $\Delta D = \SI{10}{mm}$ is the diametral
clearance. This gives $R_\mathrm{cyl} = \SI{1.605}{m}$ and a radial
gap of $\Delta r = \SI{5}{mm}$ between the outer edge of the rotating
element and the vessel wall. The vessel height is $H = \SI{1.5}{m}$.

\paragraph{Gas.}
The vessel is filled with xenon at $P = \SI{15}{bar}$ and
$T = \SI{300}{K}$ (STP). Treating Xe as an ideal monatomic gas
($M = \SI{0.131}{kg/mol}$, $\gamma = 5/3$):
\begin{alignat}{2}
  \rho  &= \frac{PM}{R_g T}
        &&\approx \SI{79}{kg/m^3}\,,
  \label{eq:rho}
  \\[4pt]
  \mu   &
        &&\approx \SI{2.32e-5}{Pa\cdot s}\,,
  \\[4pt]
  \nu   &= \mu/\rho
        &&\approx \SI{2.94e-7}{m^2/s}\,,
  \\[4pt]
  c_s   &= \sqrt{\gamma R_g T / M}
        &&\approx \SI{178}{m/s}\,.
\end{alignat}

\paragraph{Wire grid.}
A horizontal wire grid (wire diameter \SI{0.2}{mm}, centre-to-centre
spacing \SI{2}{mm}) is present in both geometries above the rotating
element. It plays a negligible role in the flow dynamics and is ignored
in the analysis that follows.


% ══════════════════════════════════════════════════════════════
\section{Geometry A: Multi-Blade Propeller}
\label{sec:propeller}
% ══════════════════════════════════════════════════════════════

\subsection{Overview}

A horizontal propeller with $N$ equally spaced blades rotates about a
vertical axis at the bottom of the cylindrical vessel. Each blade has a
NACA\,0012 airfoil cross-section at zero angle of attack ($\alpha = 0$).
The propeller sweeps through a prescribed angle and stops; the goal is
for the gas above the blades to become quiescent as quickly as possible
after stopping.

\subsection{Blade Geometry}

Each blade extends from a central hub to the tip:
\begin{itemize}
  \item Span (hub to tip): $R = \SI{1.6}{m}$.
  \item Chord: $c$ (variable parameter; default \SI{0.16}{m}).
  \item Profile: NACA\,0012 (symmetric, maximum thickness
        $t_\mathrm{max} = 0.12\,c$).
  \item Angle of attack: $\alpha = 0^\circ$.
\end{itemize}

The blade upper surface is at height $z_\mathrm{blade} = t_\mathrm{max}/2
= 0.06\,c$ above the rotation plane. For $c = \SI{0.16}{m}$, the
maximum thickness is $t_\mathrm{max} = \SI{19.2}{mm}$ and
$z_\mathrm{blade} = \SI{9.6}{mm}$.

\subsection{Propeller Layout ($N$ Blades)}

The $N$ blades are arranged at equal angular intervals of $2\pi/N$
radians ($360^\circ/N$). The total propeller diameter is $2R = \SI{3.2}{m}$,
leaving a radial tip clearance of $\Delta r = \SI{5}{mm}$ to the
vessel wall on each side.

The maximum angular sweep per event is chosen so that each blade moves
into the position vacated by its neighbour:
\begin{equation}
  \theta_\mathrm{max} = \frac{\pi}{N}
  \quad\text{(radians)}\,,
  \qquad
  = \frac{180^\circ}{N}\,.
  \label{eq:theta_max}
\end{equation}
For $N = 8$: $\theta_\mathrm{max} = \SI{22.5}{\degree}
= \SI{0.3927}{rad}$.

\subsection{Derived Geometric Quantities}

Table~\ref{tab:propA} collects the key dimensions for Geometry~A.
All quantities that depend on $c$ are given for the default
$c = \SI{0.16}{m}$ and $N = 8$.

\begin{table}[H]
\centering
\caption{Geometry A (propeller) — key dimensions ($N = 8$,
$c = \SI{0.16}{m}$).}
\label{tab:propA}
\begin{tabular}{lll}
\toprule
Quantity & Expression & Value \\
\midrule
Number of blades       & $N$                           & 8 \\
Blade span             & $R$                           & \SI{1.6}{m} \\
Blade chord            & $c$                           & \SI{0.16}{m} \\
Max blade thickness    & $0.12\,c$                     & \SI{19.2}{mm} \\
Total propeller diameter & $2R$                        & \SI{3.2}{m} \\
Vessel inner diameter  & $2R + \Delta D$               & \SI{3.21}{m} \\
Radial tip clearance   & $\Delta r = \Delta D/2$       & \SI{5}{mm} \\
Vessel height          & $H$                           & \SI{1.5}{m} \\
Max sweep angle        & $\pi / N$                     & \SI{22.5}{\degree} \\
Blade planform area (one) & $R \times c$               & \SI{0.256}{m^2} \\
Total blade area ($N$ blades) & $N R c$                & \SI{2.048}{m^2} \\
Swept area per blade   & $\tfrac{1}{2}R^2\,\theta_\mathrm{max}$
                                                        & \SI{0.503}{m^2} \\
Blockage (thickness/tip gap) & $t_\mathrm{max}/\Delta r$
                                                        & 3.84 \\
\bottomrule
\end{tabular}
\end{table}

\subsection{Boundary Conditions}

The propeller operates inside a closed cylindrical vessel:
\begin{itemize}
  \item \textbf{Bottom:} the propeller hub and mounting structure
        (solid boundary).
  \item \textbf{Top:} the vessel lid at height $H = \SI{1.5}{m}$
        above the rotation plane (solid boundary).
  \item \textbf{Side:} cylindrical vessel wall at $r = R_\mathrm{cyl}$,
        with \SI{5}{mm} radial clearance from the blade tips
        (solid boundary).
  \item \textbf{Gas:} Xe at \SI{15}{bar}, \SI{300}{K}, initially at
        rest.
\end{itemize}

The small tip clearance (\SI{5}{mm}) means the blade tips operate in a
narrow annular gap. The vessel wall constrains the flow: there is no
open outflow boundary. Any momentum deposited in the gas must either
be absorbed by the vessel walls or decay in place.

\subsection{Sequence of Operations}

For each event:
\begin{enumerate}
  \item The propeller rotates by $\theta_\mathrm{max} = \pi/N$ in a
        sweep time $T_\mathrm{sweep}$, then stops.
  \item The system waits for time $t_\mathrm{wait}$ for the gas to
        become quiescent.
  \item Measurement takes place.
\end{enumerate}

The angular velocity during the sweep, assuming constant rotation, is
$\omega = \theta_\mathrm{max} / T_\mathrm{sweep} = \pi / (N\,
T_\mathrm{sweep})$.


% ══════════════════════════════════════════════════════════════
\section{Geometry B: Rotating Disk with Carrier Plates}
\label{sec:disk}
% ══════════════════════════════════════════════════════════════

\subsection{Overview}

A solid circular disk of radius $R = \SI{1.6}{m}$ rotates about a
vertical axis at the bottom of the same cylindrical vessel. The disk
carries $N$ ion-collection plates that can slide radially on rails and
be elevated vertically above the disk surface. The disk rotation
repositions the plates; the radial slide and vertical elevation are
separate, sequential steps.

\subsection{Disk Geometry}

\begin{itemize}
  \item Radius: $R = \SI{1.6}{m}$.
  \item Thickness: $h_\mathrm{disk}$ (variable parameter; default
        \SI{5}{mm}). Must be thick enough to house the elevation
        screws for the carrier plates.
  \item The disk upper surface is at height
        $z_\mathrm{disk} = h_\mathrm{disk}$ above the vessel floor.
  \item Total disk diameter: $2R = \SI{3.2}{m}$.
  \item Radial clearance to vessel wall:
        $\Delta r = \SI{5}{mm}$ (same as propeller geometry).
\end{itemize}

The disk is a solid, flat object. Unlike the propeller blades
(which are isolated airfoils with gas flowing around them), the disk
presents a continuous solid surface. The gas above the disk interacts
with the entire upper face.

\subsection{Rail and Hangar System}

$N$ radial rails are machined into (or mounted on) the top surface of
the disk, equally spaced at angular intervals of $2\pi/N$. At the
outer perimeter of the disk, a ring-shaped \emph{hangar} houses the
$N$ carrier plates. The hangar is closed on all sides except for $N$
\emph{doors} (one per rail) through which the carriers can slide
radially inward or outward.

\subsection{Carrier Plate}

Each carrier plate is a small, flat element:
\begin{itemize}
  \item Dimensions: $\SI{160}{mm} \times \SI{160}{mm}
        \times \SI{5}{mm}$ (length $\times$ width $\times$ height).
  \item The carrier slides radially on the disk surface along its rail.
  \item It can be pushed vertically upward (by screws housed in the
        disk) so that its active surface extends above the boundary
        layer of the disk.
  \item The carrier profile must be \emph{symmetric} (top-bottom),
        since it retracts downward the same way it was elevated. A
        lenticular or symmetric diamond cross-section is appropriate.
\end{itemize}

Only \emph{one} carrier plate slides per event.

\subsection{Derived Geometric Quantities}

Table~\ref{tab:geomB} collects the key dimensions for Geometry~B.

\begin{table}[H]
\centering
\caption{Geometry B (rotating disk) — key dimensions ($N = 8$,
$h_\mathrm{disk} = \SI{5}{mm}$).}
\label{tab:geomB}
\begin{tabular}{lll}
\toprule
Quantity & Expression & Value \\
\midrule
Number of ion plates   & $N$                           & 8 \\
Disk radius            & $R$                           & \SI{1.6}{m} \\
Disk thickness         & $h_\mathrm{disk}$             & \SI{5}{mm} \\
Total disk diameter    & $2R$                          & \SI{3.2}{m} \\
Vessel inner diameter  & $2R + \Delta D$               & \SI{3.21}{m} \\
Radial tip clearance   & $\Delta r = \Delta D/2$       & \SI{5}{mm} \\
Vessel height          & $H$                           & \SI{1.5}{m} \\
Max sweep angle        & $\pi / N$                     & \SI{22.5}{\degree} \\
Disk upper surface area & $\pi R^2$                    & \SI{8.04}{m^2} \\
Carrier dimensions     & $L_c \times W_c \times h_c$   & $160\times160\times\SI{5}{mm}$ \\
Carrier planform area  & $L_c \times W_c$              & \SI{0.0256}{m^2} \\
Carrier cross-section (radial motion) & $W_c \times h_c$
                                                        & $160\times\SI{5}{mm}$ \\
Carrier cross-section (vertical motion) & $L_c \times W_c$
                                                        & $160\times\SI{160}{mm}$ \\
Elevation height       & (to clear BL)                 & TBD by analysis \\
Blockage (disk edge/gap) & $h_\mathrm{disk}/\Delta r$  & 1.0 \\
\bottomrule
\end{tabular}
\end{table}

\subsection{Boundary Conditions}

The disk operates inside the same closed cylindrical vessel as the
propeller:
\begin{itemize}
  \item \textbf{Bottom:} the vessel floor, immediately below the disk.
        The disk rests on (or is mounted just above) the floor via
        bearings.
  \item \textbf{Top:} vessel lid at $H = \SI{1.5}{m}$ above the
        vessel floor.
  \item \textbf{Side:} cylindrical wall at $r = R_\mathrm{cyl}$, with
        \SI{5}{mm} radial clearance from the disk edge.
  \item \textbf{Above the disk:} gas column of height
        $\sim H - h_\mathrm{disk} \approx \SI{1.495}{m}$.
  \item \textbf{Gas:} Xe at \SI{15}{bar}, \SI{300}{K}, initially at
        rest.
\end{itemize}

The small radial clearance (\SI{5}{mm}) means the annular gap between
the disk edge and the vessel wall is very narrow. For the default disk
thickness $h_\mathrm{disk} = \SI{5}{mm}$, this gap has an aspect ratio
of $h_\mathrm{disk}/\Delta r = 1$. Gas in this gap is confined and
strongly influenced by both the rotating disk edge and the stationary
wall.

Unlike the propeller, the disk presents a \emph{continuous solid
surface} to the gas above. The entire upper face of the disk
participates in the boundary layer, not just $N$ isolated blades.

\subsection{Sequence of Operations}

For each event, three sequential steps:
\begin{enumerate}
  \item \textbf{Disk rotation.} The disk rotates by
        $\theta_\mathrm{max} = \pi/N$ in a sweep time
        $T_\mathrm{sweep}$, then stops. The angular velocity is
        $\omega = \pi/(N\,T_\mathrm{sweep})$. During this step all
        carrier plates are housed inside the hangar; only the smooth
        disk surface is exposed to the gas.

  \item \textbf{Radial slide.} After the disk has stopped and the gas
        above the disk has become sufficiently quiescent, one carrier
        plate slides radially outward from the hangar along its rail to
        the target position. The carrier moves along the disk surface
        (i.e.\ at $z \approx h_\mathrm{disk}$). Its cross-section
        presented to the gas during this motion is
        $W_c \times h_c = 160 \times \SI{5}{mm}$.

  \item \textbf{Vertical elevation.} After the radial motion has
        stopped and any induced disturbance has decayed, the carrier is
        pushed vertically upward by a height sufficient to place its
        active surface above the disk's residual boundary layer. The
        cross-section presented to the gas during this motion is the
        full planform $L_c \times W_c = 160 \times \SI{160}{mm}$, but
        the velocity is low and the travel distance is small
        ($\sim\SI{5}{mm}$).
\end{enumerate}


% ══════════════════════════════════════════════════════════════
\section{Comparison of Geometries}
\label{sec:comparison}
% ══════════════════════════════════════════════════════════════

Table~\ref{tab:comparison} provides a side-by-side comparison of the
two geometries.

\begin{table}[H]
\centering
\caption{Side-by-side comparison of Geometry A (propeller) and
Geometry B (rotating disk). Default parameters: $N = 8$,
$c = \SI{0.16}{m}$, $h_\mathrm{disk} = \SI{5}{mm}$.}
\label{tab:comparison}
\begin{tabular}{lll}
\toprule
Property & A: Propeller & B: Disk \\
\midrule
Rotating element
  & $N$ isolated blades
  & Solid disk \\
Radius
  & \SI{1.6}{m}
  & \SI{1.6}{m} \\
Thickness
  & \SI{19.2}{mm} (blade)
  & \SI{5}{mm} (disk) \\
Cross-section
  & NACA\,0012 airfoil
  & Flat plate (sharp edge) \\
Angle of attack
  & $0^\circ$
  & N/A (flat disk) \\
Upper surface area
  & $N R c = \SI{2.05}{m^2}$
  & $\pi R^2 = \SI{8.04}{m^2}$ \\
Gas contact
  & Isolated blades with gaps
  & Continuous surface \\
Max sweep angle
  & $\pi/N = \SI{22.5}{\degree}$
  & $\pi/N = \SI{22.5}{\degree}$ \\
Radial tip clearance
  & \SI{5}{mm}
  & \SI{5}{mm} \\
Vessel height
  & \SI{1.5}{m}
  & \SI{1.5}{m} \\
BL type (during rotation)
  & Airfoil BL on each blade
  & Von K\'arm\'an rotating-disk BL \\
Edge geometry
  & Streamlined (NACA\,0012 TE)
  & Blunt (disk rim, \SI{5}{mm} step) \\
Blockage at rim
  & $t_\mathrm{max}/\Delta r = 3.8$
  & $h_\mathrm{disk}/\Delta r = 1.0$ \\
Post-stop carrier motion
  & None (blades are the sensors)
  & Radial slide + vertical elevation \\
Moving mass per event (gas disturbance)
  & All $N$ blades
  & Disk (rotation) + 1 carrier \\
\bottomrule
\end{tabular}
\end{table}

\subsection{Key Geometric Differences}

\paragraph{Surface continuity.}
The propeller exposes $N$ isolated blades to the gas; between the
blades the gas column extends unobstructed from the vessel floor to
the lid. The disk, by contrast, presents a continuous solid surface of
$\pi R^2 \approx \SI{8}{m^2}$ that separates the gas volume into two
regions: below the disk (a thin layer between disk and floor) and above
the disk (the main volume, $\sim\SI{1.5}{m}$ tall).

\paragraph{Edge geometry.}
The propeller blades have a streamlined NACA\,0012 trailing edge, which
produces a thin, low-drag wake. The disk has a blunt rim of height
$h_\mathrm{disk}$, which produces a larger separated wake at the edge.
However, the disk rim occupies only the annular region at
$r \approx R$, while the bulk of the disk surface is smooth and flat.

\paragraph{Thickness and blockage.}
The NACA\,0012 blade is much thicker than the disk
(\SI{19.2}{mm} vs.\ \SI{5}{mm} for default parameters). In the
\SI{5}{mm} radial gap, the blade blockage ratio is 3.8 (the blade is
thicker than the gap), while the disk blockage ratio is 1.0 (disk edge
just fills the gap). The propeller chord $c$ is a free parameter and
can be reduced to lower the blockage, at the cost of reduced detector
area.

\paragraph{Boundary layer character.}
On the propeller, each blade develops an independent airfoil-type
boundary layer (Blasius/turbulent flat-plate type) whose thickness
depends on chord $c$ and local Reynolds number. On the disk, the
boundary layer is a \emph{von K\'arm\'an rotating-disk} boundary layer:
it has both tangential and radial velocity components, and its thickness
scales as $\delta \sim \sqrt{\nu/\omega}$, independent of radius
(in the laminar regime). This is a qualitatively different flow.

\paragraph{Carrier motion (Geometry B only).}
In Geometry~B, after the disk stops, a carrier plate must slide
radially and then be elevated vertically. These additional motions
create their own flow disturbances, which must also decay before
measurement. Geometry~A has no such additional motion --- the blades
themselves are the sensor elements (or sit directly beneath the sensing
region).


% ══════════════════════════════════════════════════════════════
\section{Parameters for Analysis}
\label{sec:params}
% ══════════════════════════════════════════════════════════════

The flow analysis for both geometries requires the following input
parameters:

\begin{table}[H]
\centering
\caption{Input parameters for the flow analysis.}
\label{tab:params_analysis}
\begin{tabular}{llll}
\toprule
Symbol & Description & Default & Geometry \\
\midrule
$N$  & Number of blades / ion plates & 8 & Both \\
$c$  & Blade chord & \SI{0.16}{m} & A only \\
$h_\mathrm{disk}$ & Disk thickness & \SI{5}{mm} & B only \\
$T_\mathrm{sweep}$ & Sweep duration & \SI{1.0}{s} & Both \\
$z_\mathrm{obs}$ & Observation height above surface & \SI{10}{mm} & Both \\
$t_\mathrm{obs}$ & Observation time after stop & \SI{1.0}{s} & Both \\
\bottomrule
\end{tabular}
\end{table}

Fixed parameters common to both geometries are listed in
Tables~\ref{tab:propA} and~\ref{tab:geomB}, and the vessel/gas
properties in Section~\ref{sec:vessel}.

\end{document}
